% ---------------------------------------------------------------------
%  CHƯƠNG 3 — PHÂN TÍCH VÀ THIẾT KẾ
% ---------------------------------------------------------------------
\chapter{Phân tích và thiết kế}
\label{ch:thietke}

\begin{muctieu}
Chương này phân tích yêu cầu bài toán, đề xuất giải pháp, thiết kế kiến trúc
hệ thống, thiết kế dữ liệu và thiết kế mô hình gợi ý phòng trọ.
\end{muctieu}

\section{Phân tích bài toán và yêu cầu}

Hệ thống nhận đầu vào là tiêu chí tìm phòng của người dùng, gồm ngân sách
tối đa, khu vực mong muốn, diện tích tối thiểu và danh sách tiện ích bắt
buộc, và trả về danh sách phòng trọ được xếp hạng theo mức độ phù hợp kèm
điểm số. Yêu cầu chức năng gồm nhập tiêu chí tìm kiếm, tính điểm và xếp
hạng phòng trọ, hiển thị kết quả kèm bản đồ vị trí, và cho phép người dùng
lưu lại phòng trọ quan tâm để cải thiện gợi ý ở những lần tìm kiếm sau. Yêu
cầu phi chức năng gồm thời gian phản hồi dưới hai giây cho mỗi lượt tìm
kiếm trên tập dữ liệu hiện có, và giao diện dùng được trên trình duyệt phổ
thông cả ở máy tính lẫn điện thoại.

\section{Thiết kế kiến trúc hệ thống}

Kiến trúc hệ thống được trình bày ở Hình~\ref{fig:kientruc}. Dữ liệu đi
tuần tự qua bốn khối, mỗi khối là một mô-đun độc lập trong mã nguồn.

\begin{figure}[H]
\centering
\resizebox{\textwidth}{!}{%
\begin{tikzpicture}[node distance=12mm and 14mm]
  \node[objbox] (in) {Tiêu chí\\ tìm phòng};
  \node[clsbox, right=of in]  (pre) {Chuẩn hoá\\ và vec-tơ hoá};
  \node[clsbox, right=of pre] (sim) {Tính độ\\ tương đồng};
  \node[clsbox, right=of sim] (rank) {Xếp hạng\\ theo điểm số};
  \node[objbox, right=of rank] (out) {Danh sách\\ phòng gợi ý};
  \draw[flowarr] (in)   -- (pre);
  \draw[flowarr] (pre)  -- (sim);
  \draw[flowarr] (sim)  -- (rank);
  \draw[flowarr] (rank) -- (out);
\end{tikzpicture}}
\caption{Sơ đồ kiến trúc tổng thể của hệ thống.}
\label{fig:kientruc}
\end{figure}

Khối tiêu chí tìm phòng nhận đầu vào từ giao diện web. Khối chuẩn hoá và
vec-tơ hoá đưa tiêu chí người dùng và hồ sơ từng phòng trọ về cùng thang
giá trị. Khối tính độ tương đồng áp dụng Công thức~\eqref{eq:score} trên
toàn bộ phòng trọ còn trống trong cơ sở dữ liệu. Khối xếp hạng chọn ra
$k$ phòng có điểm cao nhất, đồng thời trộn thêm một số phòng trọ được người
dùng có hồ sơ tương tự đã lưu lại, trước khi trả kết quả cho giao diện.

\section{Thiết kế dữ liệu}

Bộ dữ liệu gồm 2.500 tin đăng phòng trọ tại sáu quận trung tâm Đà Nẵng,
được chia thành 2.000 phòng dùng để xây dựng hồ sơ đặc trưng phục vụ gợi ý
và 500 phòng giữ lại làm tập kiểm thử ngoại tuyến, cùng lịch sử tương tác
đã lưu hoặc đã liên hệ của 200 người dùng thử nghiệm. Cấu trúc mỗi bản ghi
phòng trọ được mô tả ở Bảng~\ref{tab:dulieu}.

\begin{table}[H]
\centering
\caption{Cấu trúc bản ghi dữ liệu phòng trọ.}
\label{tab:dulieu}
\begin{tabular*}{\textwidth}{@{\extracolsep{\fill}}>{\bfseries\color{daunavy}}l >{\color{daumaroon}}l l}
\toprule
\rowcolor{gray100}\bfseries\color{daunavy}Trường & \bfseries\color{daumaroon}Kiểu & \bfseries\color{daunavy}Mô tả \\
\midrule
id          & integer & Khoá định danh phòng trọ \\
gia\_thue   & integer & Giá thuê hằng tháng, đơn vị nghìn đồng \\
quan        & text    & Quận nơi phòng trọ toạ lạc \\
dien\_tich  & float   & Diện tích phòng, đơn vị mét vuông \\
tien\_ich   & text    & Danh sách tiện ích, phân tách bằng dấu phẩy \\
mo\_ta      & text    & Mô tả chi tiết do chủ trọ cung cấp \\
vi\_do, kinh\_do & float & Toạ độ dùng để hiển thị trên bản đồ \\
\bottomrule
\end{tabular*}
\end{table}

\section{Thiết kế mô hình gợi ý}

Mô hình được chọn là lọc dựa trên nội dung có trọng số theo
Công thức~\eqref{eq:score}, với bốn nhóm đặc trưng: giá thuê và diện tích
được chuẩn hoá theo Min-Max, khu vực được mã hoá one-hot theo quận, tiện
ích được mã hoá one-hot theo từng loại tiện ích, và mô tả văn bản được
vec-tơ hoá bằng TF-IDF với n-gram bậc một. Trọng số $w_1, \dots, w_4$ trong
Công thức~\eqref{eq:score} được chọn bằng tìm kiếm lưới trên tập kiểm thử
ngoại tuyến, tối ưu theo Precision@5, không đụng tới 500 phòng trọ giữ lại
làm tập đánh giá cuối cùng trong quá trình chọn trọng số. Thành phần cộng
tác được cài đặt đơn giản bằng cách tìm năm người dùng có vec-tơ tiêu chí
gần nhất với người dùng hiện tại, sau đó cộng thêm một phần điểm cho những
phòng trọ đã được các người dùng đó lưu lại.

\begin{luuy}
Không được dùng tập kiểm thử cuối cùng để chọn trọng số hay siêu tham số.
Làm vậy thì các độ đo báo cáo ở Chương~\ref{ch:ketqua} sẽ cao hơn thực tế,
và em sẽ không giải thích được khi hội đồng hỏi tới.
\end{luuy}
